\section{Introduzione}
\label{sec:introduction}

\subsection{Obiettivi}
Il progetto implementa \emph{da zero} i principali algoritmi di Reinforcement Learning e
li confronta su tre ambienti di difficolt\`a crescente. Come chiede la consegna,
l'enfasi \`e sulla \textbf{comprensione} e sull'\textbf{analisi sperimentale critica},
non sulla performance assoluta: diversi risultati sono volutamente sub-ottimali e
vengono analizzati come tali, inclusi i casi di fallimento e i risultati che smentiscono
le ipotesi da cui ero partito durante i primi esperimenti.

\subsection{Ambienti e algoritmi}
Ho lavorato su tre ambienti Gymnasium:
\textbf{CartPole-v1} (comune; stato 4-D continuo, 2 azioni, reward $+1$/passo, denso),
\textbf{Acrobot-v1} (assegnato; 6-D continuo, 3 azioni, reward $-1$/passo fino
all'obiettivo, sparso) e \textbf{ALE/Pong-v5} (assegnato; Atari da
pixel, tramite l'Arcade Learning Environment). La progressione non \`e casuale: CartPole serve a
verificare che gli algoritmi siano corretti, Acrobot a metterli in crisi, Pong a
mostrare cosa serve quando la tabella non \`e pi\`u praticabile.

Su CartPole e Acrobot uso quattro algoritmi tabellari scritti a mano: due di
\emph{controllo} (Q\hbox{-}Learning, SARSA) e due di \emph{predizione} (Monte Carlo,
TD(0)). Poich\'e entrambi gli ambienti hanno stato continuo, ho dovuto introdurre una
\textbf{discretizzazione}. Per Pong ho implementato un DQN in PyTorch (Mnih et al.\ 2015) con experience
replay e target network, pi\`u il bonus Double DQN (van Hasselt et al.\ 2016).

MC e TD(0), essendo di predizione, valutano una policy \emph{fissa} senza apprenderla: ho scelto di addestrare prima
Q\hbox{-}Learning, estrarne la policy greedy e valutare \emph{quella}, affiancandole una
policy casuale come riferimento (Sez.~\ref{sec:met-prediction}). La scelta non \`e
neutra: come si vedr\`a, il divario tra MC e TD(0) in termini di return dipende proprio da quale policy si
valuta.
